\section{Experimental Setup and Reproducibility}\label{app:experimental-setup}

\subsection{Models}\label{app:experimental-setup-models}

\subsubsection{Regression Models}\label{app:experimental-setup-reg-models}
Distributional random forests (DRF) are a closely related approach for estimating conditional distributions using random forests. However, to the best of our knowledge, the only publicly available DRF implementations are deprecated and not easily integrated into our experimental stack. We therefore do not include DRF as a numerical baseline, and instead compare against a broad set of actively maintained probabilistic tree and boosting methods (NGBoost, BART, CatBoost‑UQ, quantile RF, XGBoostLSS), as well as GP, Bayesian linear models, conformal methods, and deep ensembles.
\paragraph{Bayesian Linear Regression (BLR)}
Probabilistic linear model placing a prior on coefficients~\citep{bishop2006pattern}. Provides posterior distributions over predictions.

\paragraph{Bayesian Additive Regression Trees (BART)}
Ensemble of regression trees with a fully Bayesian formulation, capturing nonlinearities and uncertainty~\citep{chipman2010bart}.

\paragraph{Gaussian Processes (GP)}
Nonparametric model defining a distribution over functions; provides predictive mean and variance~\citep{williams2006gaussian}.

\paragraph{XGBoost / LightGBM / CatBoost}
High-performance gradient-boosted tree ensembles for point prediction~\citep{chen2016xgboost, ke2017lightgbm, dorogush2018catboost}.

\paragraph{NGBoost}
Probabilistic boosting framework outputting predictive distributions~\citep{duan2020ngboost}.

\paragraph{Quantile Regression Forests (QRF)}
Nonparametric method for conditional quantiles~\citep{meinshausen2006quantile}.

\paragraph{Distributional Forests (DF)}
Ensembles predicting full parametric distributions of the target~\citep{schlosser2019distributional}.

\paragraph{Conformal Prediction (MAPIE)}
Split-conformal approach to produce calibrated prediction intervals from any base model~\citep{taquet2022mapie}.

\subsubsection{Classification Models}\label{app:experimental-setup-clas-models}

\paragraph{Classical Models}
\begin{itemize}
    \item \textbf{Logistic Regression}: Linear model for binary classification estimating class probabilities via the logistic function.
    \item \textbf{Decision Trees}: Non-linear model partitioning the feature space and assigning class labels based on majority vote in each region.
    \item \textbf{Support Vector Machines (SVM)}: Finds hyperplanes separating classes, effective in high-dimensional feature spaces.
    \item \textbf{Naive Bayes}: Probabilistic classifier assuming feature independence, providing class probability estimates.
    \item \textbf{k-Nearest Neighbors (k-NN)}: Non-parametric method classifying samples based on the majority label of nearest neighbors.
\end{itemize}

\paragraph{Ensemble Methods}
\begin{itemize}
    \item \textbf{Random Forest}: Ensemble of decision trees trained on bootstrapped samples, improving accuracy through averaging.
    \item \textbf{AdaBoost}: Sequentially combines weak classifiers, focusing on misclassified instances to improve performance.
    \item \textbf{Gradient Boosting Machines (GBM)}: Sequential tree ensembles that iteratively correct previous errors.
    \item \textbf{XGBoost}: Optimized gradient boosting framework known for speed and accuracy.
    \item \textbf{LightGBM}: Histogram-based gradient boosting framework enabling faster training on large datasets.
    \item \textbf{CatBoost}: Gradient boosting framework handling categorical features efficiently without explicit encoding.
\end{itemize}

\paragraph{Probabilistic Models}
\begin{itemize}
    \item \textbf{Naive Bayes}: Provides probabilistic predictions of class labels.
    \item \textbf{Gaussian Processes (GP)}: Non-parametric model giving predictive distributions over classes.
    \item \textbf{Bayesian Additive Regression Trees (BART)}: Ensemble of trees with Bayesian inference, capturing uncertainty in predictions.
    \item \textbf{GP-BART}: Combines Gaussian Processes with BART for improved modeling of complex data structures with uncertainty.
    \item \textbf{NGBoost}: Distributional boosting. Although e.g. Exponential distribution also available always significantly worse so to ensure fair use of optuna trials we set distribution to Normal.
\end{itemize}

\paragraph{Neural Networks}
\begin{itemize}
    \item \textbf{Multi-Layer Perceptron (MLP)}: Feedforward neural network for classification.
    \item \textbf{Probabilistic Neural Networks (PNNs)}: NN with output layer being the probability of the target variable, trained via Bernoulli log-likelihood.
\end{itemize}

\paragraph{Advanced Techniques}
\begin{itemize}
    \item \textbf{NGBoost}: Gradient boosting framework that outputs probabilistic predictions.
    \item \textbf{Conformal Prediction (MAPIE)}: Provides valid prediction intervals for any underlying classifier.
\end{itemize}

\subsection{Regression Metrics}\label{app:computational-setup-regression-metrics}

\subsubsection{Point Metrics}

\paragraph{Mean Squared Error (MSE)}: Standard measure of average prediction error.
\begin{align}
    \text{MSE} = \frac{1}{N} \sum_{i=1}^N (y_i - \hat{y}_i)^2
\end{align}

\paragraph{Root Mean Squared Error (RMSE)}: Square root of MSE, less sensitive to outliers.
\begin{align}
    \text{RMSE} = \sqrt{\frac{1}{N} \sum_{i=1}^N (y_i - \hat{y}_i)^2}
\end{align}

\paragraph{Mean Absolute Error (MAE)}: Provides a robust measure of average prediction error, also less sensitive to outliers.
\begin{align}
    \text{MAE} = \frac{1}{N} \sum_{i=1}^N |y_i - \hat{y}_i|
\end{align}

\subsubsection{Probabilistic Metrics}

\paragraph{Negative Log-Likelihood (NLL)}: Evaluates the likelihood of observed targets under the predicted distribution.
\begin{align}
    \text{NLL} = -\frac{1}{N} \sum_{i=1}^N \log p(y_i \mid \mathbf{x}_i)
\end{align}

\paragraph{Continuous Ranked Probability Score (CRPS)}: Measures how closely the predicted cumulative distribution function aligns with the observed value.
\begin{align}
    \text{CRPS}(F, y) = \int_{-\infty}^{\infty} \big(F(z) - \mathbf{1}\{y \le z\}\big)^2 dz
\end{align}

\paragraph{Prediction Interval Coverage Accuracy (PICA):} As defined in KOEMEN ET AL 2025, measures the average absolute deviation between desired and empirical coverage of predictive intervals.
\begin{align}
    \text{PICA}_\alpha = |\hat{c}_\alpha - \alpha|
\end{align}

\paragraph{Sharpness}: Average width of predictive intervals; narrower intervals indicate sharper predictive distributions.

\paragraph{Interval/Weighted Interval Score (IS/WIS)}: Proper scoring rules for predictive intervals; penalizes both interval width and miscoverage.
\begin{align}
    \text{IS}_\alpha = (u-l) + \frac{2}{\alpha}(l - y) \mathbf{1}\{y < l\} + \frac{2}{\alpha}(y - u) \mathbf{1}\{y > u\}
\end{align}

\subsection{Classification Metrics}\label{app:computational-setup-classification-metrics}
\paragraph{Accuracy}: Fraction of correctly classified samples.
\paragraph{F1-Score}: Harmonic mean of precision and recall, useful for imbalanced classes.


\paragraph{Log-Loss}: Measures how well predicted probabilities align with observed class labels.
\paragraph{Brier Score}: Mean squared error between predicted probabilities and one-hot encoded labels, sensitive to calibration.
\paragraph{Expected Calibration Error (ECE)}: Measures miscalibration across probability bins.
